% Options are supplied by the test runner: stage, degree-level, lang, campus.
\documentclass{hithesis}
{\hitif[lang=en]{\bilingualswapon}{\bilingualswapoff}
\hitsetup{info={
  title={机器人宇航员双臂协调轨迹规划}[Trajectory Planning for Coordinated Motion of Dual-Armed Robotic Astronauts],
  affiliation={机电工程及自动化}[School of Mechanical Engineering],
  speciality={机械工程}[Mechanical Engineering],
  supervisor={徐教授}[Professor Xu],author={某某}[Student Name],student-id={15S1530},
  date-zh={2016-09-23},date-en={2016-09-23},
}}}
\begin{document}
\cover
\tableofcontents
\mainmatter
\ExplSyntaxOn
\typeout{SZ-REPORT:~ctex=\bool_if:NTF\g__hit_ctex_base_bool{true}{false};~word=\bool_if:NTF\g__hit_final_body_bool{true}{false};~sz=\bool_if:NTF\g__hit_shenzhen_report_bool{true}{false}}
\typeout{SZ-WIDTH:~\dim_to_decimal_in_bp:n{\textwidth}}
\typeout{SZ-HEIGHT:~\dim_to_decimal_in_bp:n{\textheight}}
\clist_map_inline:nn{Normal,Heading1,Heading2,Heading3,TOC1,TOC2,TOC3}{
  \docxlike_format_compute:n{#1}
  \typeout{SZ-STYLE:~#1;size=\dim_to_decimal_in_bp:n{\l_docxlike_format_size_dim}}
}
\ExplSyntaxOff
\hitif[lang=en]{\section{Source and objectives}}{\section{课题的来源及研究的目的和意义}}
\hitif[lang=en]{\subsection{Research background}}{\subsection{课题来源}}
\hitif[lang=en]{
Robotic astronauts can assist astronauts in constructing and maintaining space stations.
This report studies coordinated motion planning for redundant robotic arms.
The same paragraph verifies the font size, line spacing and paragraph indentation.
}{
机器人宇航员可以协助航天员完成空间站建设和维护等任务。本报告研究冗余机械臂的双臂协调运动规划。
这一段用来核对正文的字号、行距和首行缩进。正文长度足以排出数行，以便测量连续行的基线间距。
}
\hitif[lang=en]{\subsubsection{Technical approach}}{\subsubsection{研究方案}}
\hitif[lang=en]{The model uses kinematics and trajectory optimization.}{模型采用运动学建模与轨迹优化方法。}
\hitif[lang=en]{\section{Research progress}}{\section{研究进展}}
\hitif[lang=en]{This section follows the previous section on the same page.}{本节接排在上一节之后，不强制另起一页。}
\begin{enumerate}
\item \hitif[lang=en]{Build the model.}{建立模型。}
\item \hitif[lang=en]{Verify the algorithm.}{验证算法。}
\end{enumerate}
\begin{table}[htbp]
\centering
\hitif[lang=en]{\caption{Model parameters}}{\caption{模型参数}}
\begin{tabular}{cc}\toprule Parameter & Value\\\midrule $n$ & 7\\\bottomrule\end{tabular}
\end{table}
\backmatter
\begin{thebibliography}{1}
\bibitem{sample} Author. Motion planning[J]. Robotics, 2020, 1: 1--10.
\end{thebibliography}
\end{document}
